\documentclass{article}
\usepackage[utf8]{inputenc}
\usepackage[italian]{babel}
\usepackage{graphicx, subfigure}
\usepackage[absolute]{textpos}	


\title{Copertina Bicocca Unimib}

% qui c'è solo quello che serve a fare la copertina 
% non sono sicuro tutti i package siano necessari ma funzionano

\newenvironment{sistema}%
{\left\lbrace\begin{array}{@{}l@{}}}%
{\end{array}\right.}


\begin{document}
\pagestyle{empty}


\begin{textblock}{0}(1.3,0.9)
\begin{center}
\begin{tabular}{c}
{\Huge \bf Universit\`{a} degli Studi di Milano Bicocca}\\[10pt]
\hline\\[20pt]
{\bf \huge \sc Scuola di scienze} \\[10pt]
{\huge Dipartimento di Matematica e Applicazioni} \\[15pt]
\includegraphics[width=7cm]{unimib.png}\\[15pt]
{\huge Corso di Laurea Triennale in Matematica}\\[60pt]
{\hbox to 8truecm{\hrulefill}}\\[5pt]
{\bf \huge \sc Elaborato Finale}\\[15pt]
{\huge \bf TITOLO}\\[5pt]
{\hbox to 8truecm{\hrulefill}}\\[50pt]

\begin{minipage}[t]{10cm}
	{\Large{\bf Relatore:\\ Professor Proprio Lui}}
\end{minipage}\hfill\begin{minipage}[t]{5cm}\raggedleft
	{\Large{\bf Candidato: \\Tale dei Tali\\ }}
\end{minipage} \\[90pt]
{\Large Seduta del xx/yy/zzzz} \\ [5pt]

\hline\\[10pt]
\bf \LARGE Anno Accademico 202X--202Y
\end{tabular}
\end{center}
\end{textblock}
\newpage

%se volete metterlo all'interno del file con il resto della tesi mettete questi (altrimenti vi si sfasano le pagine:

%\thispagestyle{empty}
%\phantom{}
%\eject

% se volete aggiungere una pagina bianca (qui o da qualche parte) ripetete questi 3 comandi


\end{document}
